% Figure | used in parts/01-algebra/01-vectors-covectors/03-contraction.tex
%!TEX root = ../../main.tex
\begin{figure}[ht]
\centering
\begin{tikzpicture}[>=Stealth]
  \foreach \i in {0,...,7}{\draw[blue!65!black,thin] (\i*0.62,-0.45)--(\i*0.62+0.42,1.65);}
  \draw[->,red!75!black,very thick] (0.12,0.15)--(4.75,1.20);
  \node[red!75!black,below right] at (2.6,0.62) {$v$};
  \node[blue!65!black,above] at (0.15,1.65) {$\omega$};
  \node[align=left,anchor=west] at (5.9,0.65) {number of sheets pierced\\[2pt] $\Rightarrow\ \omega(v)\approx 3.9$};
\end{tikzpicture}
\caption{Geometric picture of contraction. A covector $\omega$ is a ``stack of parallel sheets'' (its level sets); the vector
$v$ is an arrow; $\omega(v)$ counts how many sheets the arrow pierces. Doubling $v$ doubles the count, and reversing $v$
changes the sign: this is linearity.}
\label{fig:sheets}
\end{figure}
